%%%%%%%%%%%%%%%%%%%%%%%%%%%%%%%%%%%%%%%%%
% Medium Length Professional CV
% LaTeX Template
% Version 3.0 (December 17, 2022)
%
% This template originates from:
% https://www.LaTeXTemplates.com
%
% Author:
% Vel (vel@latextemplates.com)
%
% Original author:
% Trey Hunner (http://www.treyhunner.com/)
%
% License:
% CC BY-NC-SA 4.0 (https://creativecommons.org/licenses/by-nc-sa/4.0/)
%
%%%%%%%%%%%%%%%%%%%%%%%%%%%%%%%%%%%%%%%%%

%----------------------------------------------------------------------------------------
%	PACKAGES AND OTHER DOCUMENT CONFIGURATIONS
%----------------------------------------------------------------------------------------

\documentclass[
	%a4paper, % Uncomment for A4 paper size (default is US letter)
	10pt, % Default font size, can use 10pt, 11pt or 12pt
]{resume} % Use the resume class

%\usepackage{ebgaramond} % Use the EB Garamond font
\usepackage{mathpazo}
\usepackage{hyperref}


\usepackage{graphicx}
\usepackage{xcolor}
\newcommand{\instlogo}[3][1.6em]{%
  \raisebox{-0.2\height}{\includegraphics[height=#1]{#2}}%
  \hspace{0.6em}\textbf{#3}%
}

%------------------------------------------------

\name{Abhijit Bendale} % Your name to appear at the top

% You can use the \address command up to 3 times for 3 different addresses or pieces of contact information
% Any new lines (\\) you use in the \address commands will be converted to symbols, so each address will appear as a single line.


\address{+1-617-697-6247 \\ abhijitbendale@gmail.com \\ US Citizen \\ Campbell, CA} % Contact information
\address{ \href{https://www.linkedin.com/in/abhijit-bendale-8175175/}{Linkedin} \quad  \\ \href{https://scholar.google.com/citations?user=C4QRs3AAAAAJ&hl=en&oi=ao}{Google Scholar} \quad  \\ \href{https://abhijitbendale.github.io/}{Website}}


%----------------------------------------------------------------------------------------

\begin{document}


%----------------------------------------------------------------------------------------
%	EXECUTIVE SUMMARY
%----------------------------------------------------------------------------------------


\begin{rSection}{Executive Summary}

Applied AI leader and hands-on engineer with 10+ years building and shipping large-scale ML systems across LLMs, multimodal AI, and computer vision. Delivered production AI platforms at global scale (400M+ users and devices).

%Career spans Founder \& CTO of Two AI, where I built multilingual LLM and multimodal AI platforms end-to-end; CTO of Samsung STAR Labs, where I led NEON, a real-time digital human platform later integrated into Samsung Electronics (\$50M+ internal valuation); and leadership roles at Samsung Electronics, where I architected Bixby Vision and on-device AI platforms deployed across 100M+ Galaxy devices.

Founder \& CTO of Two Platforms (raised 20M seed, 10M ARR), where I built multilingual LLM and multimodal AI platforms deployed at national scale. Former CTO of Samsung STAR Labs, where I led NEON, a real-time generative digital human platform later integrated into Samsung Electronics (\$50M+ internal valuation). Previously at Samsung Electronics, architected on-device AI platforms deployed across 100M+ Galaxy devices.

Deep expertise in large-scale ML systems, distributed training, post-training, RLHF, evals, Agentic AI systems. Expertise in performance optimization across model, runtime, and hardware, with a strong focus on low-latency and high-throughput AI systems. Inventor on 30+ patents with peer-reviewed papers at leading AI conferences (3,500+ citations).


\end{rSection}

%\begin{rSection}{Expertise \& Technical Skills}
%\begin{itemize}
%\item \textbf{Languages \& Core Tools}: C++, Python, TypeScript, PyTorch, CUDA / Triton, Deepspeed, Ray.
%\item \textbf{Foundation Models \& Training}: Large-scale pre/post-training of LLMs (Transformer, MoE); multimodal models (vision, audio, CLIP); alignment techniques (RLHF, DPO); efficient fine-tuning (LoRA, QLoRA); distributed training frameworks (DeepSpeed, Ray)
%\item \textbf{Inference \& Serving}: High-throughput LLM serving (vLLM, NVIDIA NIM, TorchServe); continuous batching, speculative decoding, KV-cache optimization, quantization; real-time low-latency pipelines (sub-100ms)
%\item \textbf{Retrieval \& Search Systems}: Billion-scale vector search (FAISS, Vertex Matching Engine); RAG pipelines; multimodal retrieval across text, image, audio, and video
%\item \textbf{AI Agents \& Orchestration}: Multi-agent system design, tool use, LangGraph, Google ADK
%\item \textbf{Cloud \& Infrastructure}: Production AI platforms on GCP/AWS/Kubernetes; distributed data pipelines; observability and reliability (Prometheus, Grafana)
%\end{itemize}

\begin{rSection}{Expertise \& Technical Skills}
\begin{itemize}

\item \textbf{Languages \& Core Tools}: C++, Python, TypeScript, PyTorch, CUDA/Triton, NCCL, MPI.

\item \textbf{Foundation Models \& Training}: LLM/VLM pre-training and fine-tuning; Transformers, MoE, CLIP-style multimodal models; distributed training with DeepSpeed, Ray, FSDP.

\item \textbf{Post-Training \& Alignment}: SFT, RLHF, preference optimization, reward modeling, LLM-as-judge evaluation, synthetic data, human-in-the-loop pipelines; TRL, PEFT.

\item \textbf{Inference \& Serving}: High-throughput LLM serving with vLLM, TGI, NVIDIA NIM, TorchServe; continuous batching, speculative decoding, KV-cache optimization, quantization.

\item \textbf{Retrieval \& Search Systems}: Billion-scale vector search, FAISS, Vertex Matching Engine, PgVector, RAG pipelines; multimodal retrieval across text, image, audio, and video.

\item \textbf{AI Agents \& Orchestration}: Tool use, function calling, multi-agent workflows, agent evaluation and reliability; LangGraph, LlamaIndex, Google ADK, MCP.

\item \textbf{Cloud \& Infrastructure}: Production AI platforms on GCP/AWS/Kubernetes; Vertex AI, GKE, SageMaker, EKS/ECS; distributed data pipelines, observability, reliability, and secure deployment.

\end{itemize}

\end{rSection}



% ============================================================

\begin{rSection}{Work Experience}

	\begin{rSubsection}{Two Platforms}{\textit{2021 -- 2025}}{\textbf{Founder \& Chief Technology Officer}}{Campbell, CA}
\item[] Founded Two Platforms - AI infrastructure \& applications startup (\$20M seed from Jio Platforms \& Naver Corp.) focused on multilingual LLMs, multimodal/voice AI, and enterprise applications for the APAC market. Scaled team of 20+ engineers across ML, platform and product.
\begin{itemize}

%\item Built SUTRA, multilingual LLM platform (50+ languages) with full lifecycle ownership: data, Transformer/MoE training, RLHF,  alignment, and production deployment. Built and maintained GPU cluster of 200+ GPUs for in-house training \& model serving.
%\item Developed enterprise AI stack spanning semantic search (RAG), vector embeddings, and secure enterprise deployment across cloud, VPC, and on-prem environments, contributing to \$10M+ ARR.
% \item Deployed multilingual multimodal search on Jio AI Cloud, serving 400M+ users across text, image, audio, video, and document inputs.
% %\item Built real-time voice AI and avatar systems, integrating LLM, TTS, and voice-to-voice pipelines with low-latency streaming (sub-100ms).
%\item Developed consumer AI applications like Zappy (AI chat application with 20M+ conversations) and AutoCameo (personalized video messaging platform) for national scale launch in India in partnership with Jio Platforms (IPL 2023).

\item Built SUTRA, a multilingual AI platform spanning foundation models, semantic search, voice AI. Owned the full lifecycle from data and Transformer/MoE training to RLHF, evaluation, and deployment; operated 200+ GPU infrastructure for training and serving.

\item Partnered with strategic customers including Jio Platforms to design and deploy AI applications for search, customer support, and workflow automation across cloud, VPC, and on-prem environments, contributing to \$10M+ ARR.

\item Developed enterprise AI stack including RAG, vector search, agentic workflows, and voice agents, enabling rapid deployment of production AI systems across telecom, finance, healthcare, and retail.

\item Deployed multilingual multimodal AI systems serving 400M+ users across text, image, audio, video, and documents. Built consumer products including Zappy (AI chat, 20M+ conversations) and AutoCameo (personalized video messaging), launched nationally in partnership with Jio Platforms.


\end{itemize}
	\end{rSubsection}



\begin{rSubsection}{Samsung Technology \& Advanced Research (STAR) Labs}{Sept. 2020 -- June 2021}{\textbf{Chief Technology Officer}}{Mountain View, CA}
\item[]As CTO of STAR Labs, led NEON, a real-time generative digital-human platform, from research prototype to commercialization; later integrated into Samsung Electronics at a \$50M+ internal valuation. Built and led a 40+ person multidisciplinary team across AI research, graphics, platform, and product, managing \$30M+ R\&D funding over 2 years.

\begin{itemize}
\item Conceived and led NEON, an interactive avatar platform enabling temporally stable 4K photorealistic digital humans with sub-100ms latency, launched at CES 2020.
\item Designed the core real-time multimodal generative AI stack, combining motion/audio-conditioned generation, identity personalization, neural rendering, optimized GPU inference, and streaming media delivery.
\item Built production ML/media platform spanning data pipelines, model training/evaluation, low-latency serving, observability, SDKs, cloud deployment, and enterprise integration.
\item Drove commercialization through NEON View, NEON Frame, and NEON Studio across interactive customer service, digital displays, mobile/web avatars, and enterprise media/content generation.
\item Led APAC enterprise partnerships across media, finance, retail, and telecom; supported deployments and pilots with Samsung business units and external partners.
\item Generated 10+ patents across core models, data pipelines, and system architecture; led research collaborations with MIT, Stanford, and the University of Washington.
\end{itemize}
\end{rSubsection}


	
\begin{rSubsection}{Samsung Research America}{Nov 2015 -- Aug 2019}{\textbf{Director / Principal Engineer / Staff Engineer}}{Mountain View, CA}

  \item Led architecture and scaling of on-device AI and perception platforms deployed across 100M+ Samsung Galaxy devices. Built and managed a team of 60+ engineers and researchers (20+ PhDs) with a \$15M+/year annual R\&D budget.

  \item[] \textbf{Bixby Vision \& On-Device AI}
  \begin{itemize}
  
\item Built and scaled Bixby Vision, on-device AI deployed across 100M+ Galaxy devices, enabling visual search, translation, shopping, and scene understanding.
 \item Built Samsung’s on-device AI platform from first principles, enabling ML deployment across Android, Tizen, Galaxy devices, Smart TVs, digital displays, and XR systems under strict latency, power, and memory constraints.
 \item Designed real-time camera and perception pipelines, processing live video frames for object detection, OCR, scene understanding, and AR experiences with low-latency on-device inference.
 \item Built platform SDK and APIs (Vision, AR, rendering) enabling partner applications and services to integrate with Bixby Vision, creating an extensible ecosystem.
 \item Led hardware–software co-design with Samsung Semiconductor (Exynos NPU), Qualcomm (Snapdragon), and Google (Android), optimizing inference via quantization, acceleration, and cross-platform deployment.
 \item Drove transition from applied research to production systems, defining architecture, evaluation, and deployment strategies for large-scale commercialization.

%• Developed large-scale on-device computer vision models (object recognition, OCR, scene understanding) optimized via quantization, compression, and heterogeneous execution (CPU/GPU/DSP/NPU).
%• Built 3D perception and SLAM systems for AR/XR, integrating visual, depth, and inertial sensing for real-time spatial mapping.

\end{itemize}

  \item[] \textbf{AR \& Mixed Reality Platforms}
  \begin{itemize}
\item Led AR and Mixed Reality platform initiatives in collaboration with Microsoft (HoloLens) and Google (Android / ARCore), shaping hardware–software integration and platform architecture.
\item Directed and contributed hands-on to Samsung’s in-house AR/MR stack, building SLAM and sensor fusion systems using RGB, event-based vision (DVS), IMU, and depth sensors for real-time spatial mapping.
\item Developed low-latency, power-efficient perception and rendering pipelines for XR, including patented approaches for event-driven vision and sensor fusion. \\ \\ \\
\end{itemize}
\end{rSubsection}

	 
%\begin{rSubsection}{Solariat Inc. (Acquired by Genesys Inc.)}{Dec 2012 -- Aug 2013}{Machine Learning Scientist}{San Carlos, CA}
%    \item Built production NLP pipelines for real-time sentiment classification and trend detection on Twitter scale streaming data ($>$ 10K events/sec); implemented incremental and active learning to continuously improve classifier accuracy. 
%\end{rSubsection}

\end{rSection}

%----------------------------------------------------------------------------------------
%	RESEARCH EXPERIENCE
%----------------------------------------------------------------------------------------

\begin{rSection}{Research Experience}
%------------------------------------------------
	\begin{rSubsection}
	{\textbf{SRI International}} {June 2014 -- September 2014} {PhD Intern} {Princeton, NJ}
		\item Built algorithms to automatically summarize and reconstruct sequential storylines from large-scale photo collections, modeling temporal structure across millions of images \textbf{DARPA ALADDIN}.
	\end{rSubsection}

	\begin{rSubsection}
	{\textbf{IDIAP Research Institute (EPFL)}}	{October 2013 -- January 2014}	{Visiting Student (with Prof. Sebastien Marcel)} 	{Switzerland}
		\item Researched continual learning and online model adaptation for large-scale biometric recognition under distribution shift (EU/NSF Open Vision project).
	\end{rSubsection}

	\begin{rSubsection}
	{\textbf{MIT Media Lab}}	{August 2010 -- July 2011}	{Research Scientist (with Prof. Ramesh Raskar)}	{Cambridge, MA}
		\item Developed VisionBlocks, a social computer vision platform (Google App Engine Award), and conducted image super-resolution research for DARPA ARGUS-IS.

	\end{rSubsection}

	\begin{rSubsection}
	{\textbf{MIT McGovern Institute \& Dept. of Brain \& Cognitive Sciences}} 	{July 2009 -- July 2010} 	{Visiting Student \& Research Scientist (with Prof. James DiCarlo)} 	{Cambridge, MA}
		\item Developed deep learning and biologically inspired vision models for object recognition and pedestrian detection. Research supported by \textbf{DARPA NeoVision}. 
	\end{rSubsection}

\end{rSection}

% ============================================================

%----------------------------------------------------------------------------------------
%	PUBLICATIONS
%----------------------------------------------------------------------------------------

\begin{rSection}{Publications \& Patents}
Selected list of publications (cited 3800+ times)  is listed below: 
	\begin{enumerate}
	 \item A. Bendale \textit{et.al.} ``SUTRA: Scalable Multilingual Language Model Architecture'' 2014 (\textit{arxiv 2405.06694})
	  \item A. Bendale, T. Boult ``Towards Open Set Deep Network'' CVPR 2016 (oral, 2.5K+ citations)	
	  %\item A. Bendale, T. Boult ``What do you do when you know that you don't know?" Biometrics Workshop, CVPR 2016
	  \item A. Bendale, T. Boult ``Towards Open World Recognition" CVPR 2015 (oral, 900+ citations)
   	 \item A. Bendale, T. Boult ``Reliable Posterior Probab. Estimation for Face Recognition" Biometrics W, CVPR 2014 
   	 \item A. Bendale, \textit{et.al.}, R. Raskar ``VisionBlocks: Social Computer Vision Platform" SocialCom 2011
	\end{enumerate}
	
Selected list of Patents (from 30+):  %\textit{Generating digital avatar} (\href{https://patents.google.com/patent/US11544886B2}{US 11,544,886}, 2023),
%\textit{Real-time interactive digital agents} (US App. 17/686,201, 2023),
%\textit{Interactive systems for digital humans} (US App. 17/571,106, 2023),
%\textit{CMOS-assisted inside-out DVS tracking} (\href{https://patents.google.com/patent/US11381741B2}{US 11,381,741}, 2023),
%\textit{Emotion-aware reactive interfaces} (\href{https://patents.google.com/patent/US11334376B2}{US 11,334,376}, 2023),
%\textit{Semi-dense depth estimation with DVS} (US 11,143,879, 2021),
%\textit{Vision intelligence management} (US 10,902,262, 2021),
%\textit{Semantic mapping for low-power AR} (US 10,812,711, 2020), and
%\textit{Camera pose estimation using DVS} (US 10,529,074, 2020).

\begin{enumerate}
%\item Communicative light assembly system for digital humans (\href{https://patents.google.com/patent/US12217341B2/en} US12217341B2)
\item  Generating digital avatar (\href{https://patents.google.com/patent/US11544886B2}{US 11544886}, 2023)
\item Real-time interactive digital agents (\href{https://patents.google.com/patent/US20230281466A1/en} {US  17686201}, 2023) 
%\item Interactive systems for digital humans (US App. 17/571,106, 2023),
\item CMOS-assisted inside-out DVS tracking (\href{https://patents.google.com/patent/US11381741B2}{US 11381741}, 2023)
%\item Emotion-aware reactive interfaces (\href{https://patents.google.com/patent/US11334376B2}{US 11,334,376}, 2023),
\item Semi-dense depth estimation with DVS (\href{https://patentimages.storage.googleapis.com/77/86/a3/d3c88f2dff2b7d/US11143879.pdf} {US 11143879}, 2021)
%\item Vision intelligence management (US 10,902,262, 2021),
\item Semantic mapping for low-power AR ( \href{https://patentimages.storage.googleapis.com/af/72/51/59c6812bbb8e7d/US10812711.pdf} {US 10812711}, 2020), 
%\item Camera pose estimation using DVS (US 10,529,074, 2020).
\end{enumerate}

Full list is available via \href{https://scholar.google.com/citations?user=C4QRs3AAAAAJ&hl=en&oi=ao}{Google Scholar}.

\end{rSection}

%%----------------------------------------------------------------------------------------
%%	PATENTS
%%----------------------------------------------------------------------------------------
%
%\begin{rSection}{Patents}
%\textit{Inventor on 30+ granted and pending patents spanning on-device AI, AR/VR, SLAM, novel sensors, digital humans, and computer vision. Full patent list available via }
%\href{https://scholar.google.com/citations?user=C4QRs3AAAAAJ&hl=en&oi=ao}{Google Scholar}.
%
%\begin{enumerate}
%\item ``Generating digital avatar'' 2023. US Patent 11,544,886
%\item ``Autonomous generation, deployment, and personalization of real-time interactive digital agents'' 2023. US Patent App. 17/686,201
%\item ``Interactive system application for digital humans'' 2023. US Patent App. 17/571,106
%\item ``CMOS-assisted inside-out dynamic vision sensor tracking for low power mobile platforms'' 2023. US Patent 11,381,741
%\item ``Emotion-aware reactive interfaces'' 2023. US Patent 11,334,376
%\item ``Semi-dense depth estimation from a dynamic vision sensor (DVS) stereo pair and a pulsed speckle pattern projector'' 2021. US Patent 11,143,879
%\item ``Vision intelligence management for electronic devices'' 2021. US Patent 10,902,262
%\item ``Semantic mapping for low-power augmented reality using dynamic vision sensor'' 2020. US Patent 10,812,711
%\item ``Camera pose and plane estimation using active markers and a dynamic vision sensor'' 2020. US Patent 10,529,074
%\end{enumerate}
%
%\end{rSection}


%----------------------------------------------------------------------------------------
%	EDUCATION SECTION
%----------------------------------------------------------------------------------------

\begin{rSection}{Education}
	
	\textbf{University of Colorado} \hfill \textit{2011 -- 2015} \\ 
	Ph.D. in Computer Science (Dissertation: Open World Recognition) \\
	Advisor: Prof. Terrance Boult \\
	\textbf{University of Colorado} \hfill \textit{2007 -- 2009} \\ 
	M.S. in Computer Science \\
	\textbf{University of Pune} \hfill \textit{2003 -- 2007} \\ 
	B.E. in Electronics and Telecommunications 
	
\end{rSection}


%----------------------------------------------------------------------------------------
%	HONORS & AWARDS
%----------------------------------------------------------------------------------------

%\begin{rSection}{Honors \& Awards (Selected)}
%
%	\begin{enumerate}
%	\item CVPR 2015 Doctoral Consortium Award (one of 35 top PhD students in Computer Vision)
%	\item IEEE BTAS 2015 Doctoral Consortium Award  (one of 10 top PhD students in Biometrics)
%	\item NSF Data Science 2015 Invite (one of top 100 PhD students in STEM)
%	\item NVIDIA 2013 CUDA Teaching Center of Excellence for University of Colorado
%	\item Google App Award 2011 for Vision Blocks project at MIT Media Lab
%	\end{enumerate}
%
%\end{rSection}

%----------------------------------------------------------------------------------------
%	TECHNICAL STRENGTHS SECTION
%----------------------------------------------------------------------------------------

%\begin{rSection}{Technical Skills \& Activities}
%
%	\begin{tabular}{@{} >{\bfseries}l @{\hspace{6ex}} l @{}}
%		Reviewing & CVPR, ICCV, NeurIPS, ICLR, WACV, IJCB, IEEE Trans. Affective Comp., NeuroComputing, IJCV \\
%		Teaching &  CS2060 Programming with C,  CS4440/5440 Big Data  \\
%		Languages & C, C++, Python, Swift, Shell, CUDA, Objective C etc. \\
%		Cloud Platforms & AWS, GCP, Azure (Cloud Deployment, High Performance Computing etc.) \\
%		Research Interests & CV, ML, Generative AI, LLMs, Image / Video Generation, Neural Rendering, SLAM, On-Device AI  \\
%		
%	\end{tabular}
%
%\end{rSection}

%----------------------------------------------------------------------------------------
%	EXAMPLE SECTION
%----------------------------------------------------------------------------------------

%\begin{rSection}{Section Name}

	%Section content\ldots

%\end{rSection}

%----------------------------------------------------------------------------------------

\end{document}
